\section{Regras de validação}\label{sec:regras}

Esta secção resume as regras que a aplicação interativa usa para descodificar e
validar uma matrícula. São as regras descritas neste documento, escritas de forma
a poderem ser verificadas. Cada problema encontrado tem um de três graus:

\begin{description}
  \item[Impossível] a matrícula não pode existir segundo as regras conhecidas;
  \item[Anacronismo] a matrícula pode existir, mas não na época que a série sugere;
  \item[Por confirmar] depende de uma regra ou entrada de tabela que não foi
        possível confirmar numa segunda fonte.
\end{description}

Os separadores (traços, espaços, pontos) não são obrigatórios. Uma matrícula
escrita com os separadores do formato é reconhecida como \emph{exata}; sem eles,
como \emph{possível}. Algumas matrículas têm mais de uma leitura (por exemplo,
\plate{AP-12-34} é um automóvel ou um reboque da Marinha) e todas são apresentadas.

\subsection{Série de 1911}
\begin{itemize}
  \item Formato \plate{Z-0} a \plate{Z-00000}, com \plate{Z} = \plate{N}, \plate{C},
        \plate{S} (1911–1936), \plate{A} ou \plate{M} (1918–1962).
  \item Sufixo \plate{-A}: aluguer; sufixo \plate{-WW}: provisória.
  \item \plate{MG-000}: Exército.
  \item Conversão para 1937 conforme a secção 1937. Para os Açores, a Madeira,
        Coimbra a partir de \plate{C-10000} e Lisboa a partir de \plate{S-40000} a
        fonte não indica a série de destino.
\end{itemize}

\subsection{Série de 1937 (\plate{XX-00-00})}
\begin{itemize}
  \item Zona pelo grupo de letras (até ao início dos anos 80) e grupos especiais
        das ilhas, serviços regionais, importação temporária, antigas colónias e
        forças militares (anexo~\ref{sec:grupos}).
  \item \plate{Q}, \plate{J}, \plate{CU} e \plate{FD}: anacronismo; só foram
        usados a partir do início dos anos 80.
  \item \plate{K}, \plate{W} e \plate{Y}: impossível.
\end{itemize}

\subsection{Série de 1992 (\plate{00-00-XX})}
\begin{itemize}
  \item \plate{KA}–\plate{KF}: importação em segunda mão (1997);
        \plate{KG}: importação por particular (1997–1998, por confirmar);
        outros grupos \plate{K}: por confirmar.
  \item Grupos começados por \plate{W} ou \plate{Y}: impossível (reservados a
        importações, nunca usados).
  \item \plate{K}, \plate{W} ou \plate{Y} na segunda letra: impossível.
\end{itemize}

\subsection{Série de 2005 (\plate{00-XX-00})}
\begin{itemize}
  \item \plate{K}, \plate{W} e \plate{Y}: impossível.
\end{itemize}

\subsection{Série de 2020 (\plate{XX 00 XX})}
\begin{itemize}
  \item Vogal na 2.ª ou na 6.ª posição: impossível, exceto se um dos grupos for
        uma vogal repetida (\plate{AA}, \plate{EE}, \plate{II}, \plate{OO},
        \plate{UU}).
  \item \plate{AA 00 AA}: impossível (reservada, nunca emitida).
  \item \plate{K}, \plate{W} e \plate{Y}: por confirmar.
  \item A ordem de atribuição presumida é: algarismos, depois o grupo da direita,
        depois o da esquerda. É coerente com a primeira matrícula
        (\plate{AA 01 AA}), mas o IMT não a publicou.
\end{itemize}

\subsection{Outros formatos}
\begin{itemize}
  \item \textbf{Ciclomotores municipais} \plate{0-XXX 00-00}: \plate{XXX} tem de
        constar da tabela de municípios; código desconhecido: impossível.
  \item \textbf{Reboques civis}: código da tabela de serviços distritais; com um
        código de uma letra o número tem até 6 algarismos, com duas letras até 5.
        Em Lisboa, a data é estimada pela tabela da secção de reboques.
        Bragança tem dois códigos nas fontes (\plate{BG} e \plate{BN}); ambos
        ficam por confirmar.
  \item \textbf{Reboques militares}: \plate{0000-M}, \plate{E-00-00}, \plate{AP-00-00}.
  \item \textbf{GNR} \plate{GNR X-0000}: função desconhecida: impossível.
  \item \textbf{Exportação} \plate{00000-X}, com \plate{X} = \plate{L}, \plate{P},
        \plate{A} ou \plate{M}; um número de algarismos diferente de 5 fica por
        confirmar.
  \item \textbf{Diplomáticas} \plate{000-CD000}, \plate{000-CC000},
        \plate{000-FM000}: código de país fora da tabela fica por confirmar.
\end{itemize}

\subsection{Palavras nas matrículas}

A aplicação procura palavras de seis letras que se podem ler numa matrícula
\plate{XX 00 XX}, lendo os algarismos do meio como letras
(\plate{0}=O, \plate{1}=I, \plate{2}=Z, \plate{3}=E, \plate{4}=A, \plate{5}=S,
\plate{6}=G, \plate{7}=T, \plate{8}=B, \plate{9}=P). Por exemplo, \plate{AB 47 EI}
lê-se “ABATEI”. Na lista de 13\,046 palavras do projeto há 3\,321 matrículas
destas, mas só 510 respeitam a regra das vogais e podem ser emitidas.
